%% ============================================================================
%% Section: Incremental Index Maintenance
%% Follows CODE-PLAN.md WP3 (patch, cone recompute, flips, validation) and
%% WP4 (derived worlds by graph surgery).
%% ============================================================================
\section{Incremental Index Maintenance}\label{sec:incr}

A curated causal graph is edited continually: an edge is added as evidence
accrues, a variable is reclassified as latent, an edge is deactivated pending
review. Requirement R4 asks that the compiled index survive these edits
without a full recompilation after every change. In this section we describe
how Causal4j observes graph mutations, validates them, and patches the
snapshot within the region of the graph they affect.

\noindent\textbf{Observing and validating changes.} The plugin registers a
transaction-event listener, scoped to the database for its lifetime, that
sees every committing write. Before a transaction commits, the listener
inspects the transaction data and, if it touches the causal schema, distills
the change into a compact change batch: created and deleted \textsf{Variable}
nodes, created and deleted \textsf{CAUSES} edges, edges whose \textsf{active}
flag flipped (treated as insertions or deletions), and variables whose
\textsf{latent} flag flipped. Transactions that touch no \textsf{Variable}
node and no \textsf{CAUSES} edge are rejected by a guard on the transaction
data and cost nothing. The pre-commit hook is also where the DAG invariant is
enforced: for every inserted edge $u \rightarrow v$, the listener tests
whether $v$ reaches $u$ in the current closures or through the other edges of
the same batch, and aborts the transaction if so, so that an invalid model
never reaches the store. The batch is handed to the post-commit hook, which
applies it once the write is durable.

\noindent\textbf{Insertion patch.} Given a change batch, the closures of
Section~\ref{sec:index:bitsets} can be maintained locally rather than
rebuilt. Inserting an edge $u \rightarrow v$ extends the CSR adjacency and
propagates closures over the affected cone only: every node in
$\an(u) \cup \{u\}$ gains $\textsf{des}[v] \cup \{v\}$ as new descendants, and
every node in $\de(v) \cup \{v\}$ gains $\textsf{anc}[u] \cup \{u\}$ as new
ancestors. Both gains are one bitset each, read from the snapshot before the
patch, and each affected row receives one word-parallel union, so the patch
costs $(|\an(u)| + |\de(v)| + 2)\,w$ word operations. Layers are relabeled
for $v$ and its descendants only when $\textsf{layer}[v] \le
\textsf{layer}[u]$. Nodes outside the two cones are untouched.

\noindent\textbf{Deletion by cone recompute.} Deleting $u \rightarrow v$ can
shrink closures in a way that no local union expresses, since a node may
have lost one of several paths. However, the rows that can change are still
bounded: only the ancestor sets of $C = \de(v) \cup \{v\}$ and the descendant
sets of $C' = \an(u) \cup \{u\}$ pass through the deleted edge. We therefore
recompute $\textsf{anc}$ over $C$ in topological order from the updated
parents, and $\textsf{des}$ over $C'$ in reverse topological order from the
updated children, with the recurrence of Section~\ref{sec:index:bitsets}
restricted to those rows; rows outside $C \cup C'$ are provably unchanged.
The cost is $O((|C| + |C'| + |E_{C \cup C'}|)\,w)$, and full recompilation
remains the fallback when $|C| + |C'|$ exceeds a threshold, which
Section~\ref{sec:eval} measures. Adding a variable allocates a dense index
and, one time in 64, one more word per row; removing a variable deletes its
incident edges and frees its index. A \textsf{latent} flip is a property-only
update that leaves the closures untouched.

\noindent\textbf{Copy-on-write and consistency.} A patch never mutates the
snapshot that readers hold. It produces a new snapshot that shares every
untouched row and array with the old one and replaces the rows of the
affected cones, and it is published with a single reference store, exactly
as a recompile is. In this way the snapshot-isolated consistency model of
Section~\ref{sec:compile} is preserved without locks, and the memory cost of a
patch is proportional to the cone rather than to the graph. Batches are
applied by first recomputing the union of the deletion cones once and then
applying the insertion unions, so that a transaction with many edits costs
one cone recompute rather than one per edit.

\subsection{Derived worlds by graph surgery}\label{sec:incr:worlds}

An intervention $do(X)$ defines a derived model, the same DAG with the edges
into $X$ removed, and the language of Section~\ref{sec:lang:registry} lets an
analyst name such a world and query it. The deletion machinery above is
exactly what constructs it: \texttt{causal.intervene} applies a batch that
deletes the incoming edges of every intervened variable to a copy of the base
snapshot, recomputes the affected cones, and registers the result under the
world's name. Since the copy shares every untouched row with the base, a
derived world costs the cone recompute in time and the changed rows in
memory, and a later structural statement over it reads its own closures like
any other snapshot. The same mechanism gives \textit{what-if} structural
queries over hypothetical edits (a proposed edge, a removed confounder)
without touching the stored graph.
